% Generated draft: requires mathematical and source review.
\documentclass[11pt]{article}
\usepackage[margin=25mm]{geometry}
\usepackage{fontspec}
\setmainfont{Latin Modern Roman}
\usepackage{amsmath,amssymb,mathtools,unicode-math}
\setmathfont{Latin Modern Math}
\usepackage{enumitem,needspace,xurl,hyperref}
\setlength{\parindent}{0pt}
\setlength{\parskip}{5pt}
\setlength{\emergencystretch}{3em}
\newcommand{\N}{\mathbb{N}}
\newcommand{\Z}{\mathbb{Z}}
\newcommand{\Q}{\mathbb{Q}}
\newcommand{\R}{\mathbb{R}}
\newcommand{\C}{\mathbb{C}}
\newcommand{\F}{\mathbb{F}}
\newcommand{\A}{\mathbb{A}}
\newcommand{\PP}{\mathbb{P}}
\newcommand{\E}{\mathbb{E}}
\newcommand{\eps}{\varepsilon}
\newcommand{\dd}{\,\mathrm{d}}
\begin{document}
\section*{UnsolvedMath ID: 20001301}
\textbf{Source problem code:} catalogue entry
\subsection{Definitions and mathematical statement}
This record accidentally joins two consecutive problems, numbered 3 and 4. Their mathematical targets and supplied status comments are kept separate.

\textbf{Problem 3: radial dualization.} Let $\mathcal K^n$ be the class of convex bodies in $\R^n$, with $n\geq2$. For $a,b\geq0$ the Minkowski combination is $aK+bL=\{ax+by:x\in K,y\in L\}$. A Minkowski class $\mathcal M$ is closed in Hausdorff distance and under Minkowski combinations and translations. A generalized $\mathcal M$-body is a convex body $K$ with $K+M_1=M_2$ for some $M_1,M_2\in\mathcal M$. For a subgroup $G\leq\operatorname{GL}(n)$, the class is $G$-invariant if $gK\in\mathcal M$ whenever $K\in\mathcal M$ and $g\in G$. Let $\mathcal M_{B,G}$ be the smallest such invariant Minkowski class containing $B$.

The source attributes to Schneider and F. Schuster the theorem: for every non-symmetric $B\in\mathcal K^n$ and every Hausdorff neighborhood $U$ of $B$, there is an affine image $B'\in U$ of $B$ such that generalized $\mathcal M_{B',\operatorname{SO}(n)}$-bodies are dense in $\mathcal K^n$. Here $\operatorname{SO}(n)$ is the rotation group. The source does not specify whether ``non-symmetric'' excludes symmetry about any center or only about the origin.

R. Gardner asks for a dual theorem in which Minkowski sums are replaced by radial sums and generalized zonoids by ``generalized intersection bodies''. For a star body with continuous positive radial function $\rho_K$ about $0$, radial sum and positive scaling are described by $\rho_{K\widetilde+L}=\rho_K+\rho_L$ and $\rho_{aK}=a\rho_K$. A possible precise rotation-generated model is the closed cone generated by $\{\rho_B\circ R^{-1}:R\in\operatorname{SO}(n)\}$ in the uniform radial metric $d_r(K,L)=\|\rho_K-\rho_L\|_\infty$; a generalized radial body has a positive radial function expressible as a difference of two members of this cone. This is the model addressed by the supplied research aid, not an asserted unique definition of the source's proposed duality. The dual of translation closure, the permitted affine transformations, target class, symmetry, and the meaning of generalized intersection bodies remain unspecified in the original question.

In that model, the aid reports that generalized rotation-generated radial bodies are dense in all star bodies exactly when the projection of $\rho_B$ to every spherical-harmonic degree is nonzero. For origin-symmetric generators and targets the criterion uses every even degree instead. It also reports that standard intersection bodies have even radial functions and hence cannot be dense among nonsymmetric star bodies, whereas generic arbitrarily small spheroidal perturbations of the ball satisfy the even-degree criterion and generate density in the symmetric category. These are attributed partial formulations with the uniform radial topology, not a complete resolution of every possible dualization.

\textbf{Problem 4: segment complexity of zonoid approximation.} A zonotope is a Minkowski sum of finitely many line segments; a zonoid is a Hausdorff limit of zonotopes. Given an $n$-dimensional zonoid $K$ and tolerance $\varepsilon>0$, the question is how many segments suffice for a zonotope $Z$ with
\[
 d(Z,K)\leq1+\varepsilon.
\]
The source says Talagrand's bound has the form $M\leq C(\varepsilon)n\log n$, while for the Euclidean ball $C(\varepsilon)n$ suffices, and asks whether the extra $\log n$ is necessary for general zonoids. Its distance $d$ is not defined. If the intended distance is Banach--Mazur distance, then, after translating centrally symmetric bodies to the origin, this means
\[
 d_{\mathrm{BM}}(Z,K)=\inf\{\lambda\geq1:\exists T\in\operatorname{GL}(n),\ Z\subseteq TK\subseteq\lambda Z\}.
\]
This interpretation requires source confirmation. For each fixed tolerance the requested improvement is a segment bound $M\leq C(\varepsilon)n$ uniformly over $n$ and $K$; the assertion is an existence upper bound, not that every good approximating zonotope must have that many segments.

The supplied research summary attributes removal of the logarithmic dimension factor to Reis and Rothvoss, arXiv:2606.28147 (June 2026), reporting $O(n\varepsilon^{-2}\log(1/\varepsilon))$ segments in the small-tolerance regime. No exact tolerance range, metric convention, or full theorem is quoted in the record. The aid also reports an elementary $O_{r,\varepsilon}(n)$ approximation result for direct sums of zonoid blocks of dimension at most a fixed $r$. These are supplied status and partial-result statements; no independent literature verification is claimed.
\subsection{Short English statement}
Two questions were fused: whether a rotation-generated Minkowski density theorem has a radial/intersection-body analogue, and whether every zonoid admits a linear-size zonotope approximation at fixed accuracy. The supplied aid gives a radial harmonic criterion under a particular model and reports a June 2026 removal of the dimension logarithm for the second question. Missing duality definitions and the approximation distance are left explicit.
\subsection{Sources}
\begin{itemize}
\item Local UnsolvedMath registry, permanent ID 20001301, exact source code catalogue entry; canonical \texttt{aim-convex-geometry-notes.json}, \texttt{notes[32]}, source item 3, tag ``problem'', workshop \emph{Fourier analytic methods in convex geometry}. Supplied URL: \url{https://aimath.org/WWN/fourierconvex/fourierconvex.pdf}. The verification explicitly identifies an accidental merger of Problems 3 and 4 but truncates before its announced separation.
\item Original attributions: Schneider and F. Schuster for the Minkowski density theorem; R. Gardner for the dualization question; Talagrand for the $C(\varepsilon)n\log n$ zonotope bound. No full bibliographic entries or separate proposer for Problem 4 are supplied.
\item The research summary cites Reis and Rothvoss, arXiv:2606.28147, June 2026, for removal of Talagrand's $\log n$ factor. Only that author/identifier/date information and the asymptotic bound are supplied; no paper title is invented.
\item Local metadata: PARTIAL-PROGRESS, registry status ``partially\_solved'', statement status ``reconstructed\_unverified''; attempt 1, result ``partial\_result'', problem status at run ``partially\_solved'', recorded rigor ``proved'', novelty confidence ``low'', candidate type ``theorem''. Original remarks and literature field are empty. The aid is machine generated, not peer reviewed, and must be independently verified before citation. No outside source was visited.
\end{itemize}
\end{document}
